\section{Small stabilizers}\label{sm:section}

We construct homomorphisms whose values belong to the local error
ideals, but are large enough to absorb a prescribed countable family
of errors.  All matrix comparisons in this section are stable
comparisons.  The identifications of
\Cref{st:folding} are understood; in particular, a rational UHF
fraction changes ranks, not operator norms.

\begin{theorem}[Small stabilizers]\label{sm:stabilizers}
Let \(S=C_0(Y)\otimes P\), and prescribe countably many pairs
\[
 x_l\in S_+,
 \qquad e_l\in J(U_{x_l}),
 \qquad l\in\N.
\]
There is a homomorphism \(\theta:S\to D\) satisfying
\begin{equation}\label{sm:local-and-lower}
 \theta(S(U))\subseteq J(U),
 \qquad
 e_l\in\Ideal_D(\theta(x_l))
 \quad(U\subseteq X\text{ open},\ l\in\N),
\end{equation}
in either of the following situations.
\begin{enumerate}[label=\textup{(\roman*)}]
\item An existing model \(\chi:S\to D\) is given.
In this case \(\theta\) can be a slowly vanishing rational UHF
fraction of \(\chi\), and it also satisfies
\begin{equation}\label{sm:coarse-fullness}
 \theta(S(U_x))\subseteq\Ideal_D(\theta(x))
 \qquad(x\in S_+).
\end{equation}
For every fixed \(N\), \(N\) copies of its defining fraction
fit in the defining copy of \(\chi\) at all sufficiently late
coordinates.
\item \(Y=(0,1]\).  No model is needed in this case, and
\eqref{sm:coarse-fullness} is not asserted.
\end{enumerate}
In both cases, for every \(l\) and \(\eps>0\) there is
\(\delta_{l,\eps}>0\) such that
\begin{equation}\label{sm:fixed-lower-cut}
 [ (|e_l|-\eps)_+]
 \leq[(\theta(x_l)-\delta_{l,\eps})_+]
 \quad\text{in }\Cu(D).
\end{equation}
The construction permits the prescribed errors and the stabilizer
to occupy separate matrix corners.
\end{theorem}

The cut in \eqref{sm:fixed-lower-cut} will remain fixed while the
accuracy, matrix sizes, and finite lists change.  We first describe
the countable tests and the comparison convention used to ensure
this.

\subsection{Fixed tests and comparison bounds}

Choose a countable compact-containment basis for the open ideals
as in \Cref{st:controllers}.  For every pair \(V\Subset U\) in this
basis, choose a norm-dense sequence of positive contractions in
\(S(V)\), and include all positive rational cut tolerances.
Enumerate the resulting upper tests as
\begin{equation}\label{sm:upper-tests}
 (a_i,\eps_i,V_i,U_i),\qquad
 a_i\in S(V_i)_+,\quad V_i\Subset U_i,
 \quad \eps_i>0.
\end{equation}
For each test fix a compactly Cuntz-supported positive
\(b_i\in Q(U_i)\) whose generated ideal contains \(Q(V_i)\).
These controllers are chosen once, before the approximation
stages.  Showing
\begin{equation}\label{sm:upper-goal}
 d_\tau((\theta_n(a_i)-\eps_i)_+)
 \leq s_{i,n}d_\tau(b_i),
 \qquad s_{i,n}\longrightarrow0,
\end{equation}
uniformly in \(\tau\), suffices for local smallness.
Indeed, \Cref{seq:smallideal-properties} first gives the conclusion
on the dense positive tests in \(S(V_i)\).  Linearity and norm
closure then give it on \(S(V_i)\), and the ideals with
\(V_i\Subset U\) exhaust \(S(U)\).

Choose bounded representatives \(e_{l,n}\) for the errors.
Enumerate every error together with a sequence of positive cut
tolerances decreasing to zero.  Write a lower request as
\(q=(l,\eta)\), and put
\begin{equation}\label{sm:error-controllers}
 A_{q,n}=(|e_{l,n}|-\eta)_+,
 \qquad
 d_\tau(A_{q,n})\leq r_q(n)d_\tau(c_q),
 \qquad r_q(n)>0,\quad r_q(n)\longrightarrow0.
\end{equation}
Here \(c_q\) is a fixed compactly Cuntz-supported positive
in \(Q(U_{x_l})\), and the inequality holds eventually for
every \(\tau\in T(Q)\).  Enlarging finitely many values of
\(r_q\) will not matter.  A request with zero controller has
eventually zero source cut and can be omitted.

We will use the fixed-cut comparison bound from
\Cref{seq:uniform-witness}.  Its elementary form is worth recording.
If \(A\precsim(B-\delta)_+\), choose \(v\) with
\(v^*(B-\delta)_+v\) close to \(A\), and set
\[
 w=g_\delta(B)v,
 \qquad
 g_\delta(t)=\sqrt{\frac{(t-\delta)_+}{t}},
 \qquad g_\delta(0)=0.
\]
Then
\begin{equation}\label{sm:witness-bound}
 w^*Bw=v^*(B-\delta)_+v,
 \qquad
 \|w\|^2\leq
 \frac{\|A\|+\|v^*(B-\delta)_+v-A\|}{\delta}.
\end{equation}
Thus coordinate comparisons into a fixed positive cut have
bounded witnesses into the uncut target.  This conclusion is
independent of matrix size.

\subsection{Fractions of an existing model}

\begin{proof}[Proof of \Cref{sm:stabilizers}\textup{(i)}]
Take c.p.c. lifts \(\chi_n:S\to Q\) of \(\chi\), using
nuclearity and the Choi--Effros lifting theorem
\cite[Theorem~3.10]{ChoiEffrosLifting}.
The constant-ideal assertion of \Cref{seq:all-traces} gives
\[
 \Ideal_D(\chi(x_l))=\Ideal_D(Q(U_{x_l})).
\]
Compact support of the controller in
\eqref{sm:error-controllers}, with an extra cut before passing
to coordinates, consequently gives an integer \(M_q\) and
\(d_q>0\) such that
\begin{equation}\label{sm:model-lower}
 d_\tau(c_q)
 \leq M_q d_\tau((\chi_n(x_l)-d_q)_+)
 \quad\text{eventually, for every }\tau.
\end{equation}
Here is the cut detail.  Write
\([c_q]\leq[(z-\gamma)_+]\), with \(z\in Q(U_{x_l})_+\).
The constant \((z-\gamma/2)_+\) lies in the ideal generated
by \(\chi(x_l)\).  Its suitable smaller cut is therefore
below finitely many copies of a positive cut of \(\chi(x_l)\).
Leave a further cut between \((z-\gamma)_+\) and that source
before lifting the comparison.  The source perturbation lemma
then gives \eqref{sm:model-lower} at every late coordinate.
Both its multiplicity and target tolerance depend only on the
request, not on the coordinate.

For an upper test, choose a strictly positive $h_i\in Q(V_i)$.
The constant-ideal assertion gives
$\chi(a_i)\in\Ideal_D(Q(V_i))=\Ideal_D(h_i)$.
Compact containment between two positive cuts of $\chi(a_i)$
therefore gives a finite integer $N_i$ and $\kappa_i>0$ with
\[
 [(\chi(a_i)-\eps_i/2)_+]\leq N_i[(h_i-\kappa_i)_+].
\]
Let $c_i\in Q(V_i)_+$ be the folded sum of these $N_i$ cuts.
It is compactly Cuntz supported in $Q(V_i)$.
The common-controller lemma then gives $[c_i]\leq L_i[b_i]$
for a finite integer $L_i$.  Lift a comparison for the less deeply
cut source with error smaller than the remaining source slack.
The positive perturbation lemma gives
\begin{equation}\label{sm:model-upper}
 d_\tau((\chi_n(a_i)-\eps_i)_+)
 \leq L_i d_\tau(b_i)
 \quad\text{eventually, for every }\tau.
\end{equation}
Thus both the controller $b_i$ and its multiplicity are fixed
before the UHF fractions are selected.

At stage \(j\), choose a positive rational \(t_j\) with
\[
 t_j\leq j^{-1},
 \qquad t_jL_i\leq j^{-1}\quad(i\leq j).
\]
Only after choosing \(t_j\), choose an increasing sequence
of coordinate thresholds \(N_j\) so that all the first
\(j\) coordinate comparisons hold and
\begin{equation}\label{sm:model-delay}
 M_qr_q(n)\leq t_j
 \qquad(q\leq j, n\geq N_j).
\end{equation}
For \(N_j\leq n<N_{j+1}\), let \(p_n\) be a projection
of normalized trace \(t_j\) in an extra UHF factor and use
\[
 \theta_n(a)=\chi_n(a)\otimes p_n.
\]
After the fixed trace-preserving identifications these are
c.p.c. maps into \(Q\).  Their multiplicative defects tend
to zero because those of \(\chi_n\) do, so they define a
homomorphism \(\theta:S\to D\).
Equation~\eqref{sm:model-upper} gives
\eqref{sm:upper-goal} with \(s_{i,n}\leq j^{-1}\) on the
\(j\)-th stage.  Moreover,
\begin{align}
 d_\tau(A_{q,n})
 &\leq r_q(n)d_\tau(c_q)\notag\\
 &\leq M_qr_q(n)d_\tau((\chi_n(x_l)-d_q)_+)
 \leq d_\tau((\theta_n(x_l)-d_q)_+).
 \label{sm:fraction-lower}
\end{align}
All inequalities use extended nonnegative arithmetic;
the scalar factors are strictly positive.  Pointwise order
in \(\Cu(Q)\) gives the coordinate comparison.  Apply
\eqref{sm:witness-bound} to the target
\((\theta_n(x_l)-d_q/2)_+\), cut further at \(d_q/2\).
Bounded witnesses now pass to \(D\), proving
\eqref{sm:fixed-lower-cut} for the enumerated tolerances,
and hence for every tolerance.

It remains to check \eqref{sm:coarse-fullness}.  If
\(y\in S(U_x)_+\), then
\(\chi(y)\in\Ideal_D(\chi(x))\).  For every source
tolerance, a slightly larger source cut is below finitely
many copies of a fixed cut of \(\chi(x)\).  Passing to
coordinates with slack gives this rank inequality eventually.
Tensoring both sides with the same \(p_n\) multiplies both
ranks by its positive trace, so the inequality remains true.
Pointwise coordinate order and
\eqref{sm:witness-bound} give
\((\theta(y)-\eps)_+\in\Ideal_D(\theta(x))\).
Letting \(\eps\downarrow0\), then using linearity and
closure, proves \eqref{sm:coarse-fullness}.
Finally \(Nt_j\leq1\) eventually for every fixed \(N\),
so \(N\) equal rational subprojections fit in the original
UHF unit.  They give the stated copies inside the spare model.
\end{proof}

\subsection{Cone blocks with a fixed lower signal}

Without a model to cut down, we must first build finite matrix maps
whose values detect each requested controller.  The construction has
three ingredients: a fixed quotient-norm threshold for the controlling
input, finitely many positive support elements with plateaux, and
finite-dimensional approximations of the corresponding quotient cones.
The plateaux convert a matrix norm lower bound into a Cuntz comparison
at the fixed threshold.  Stage-dependent matrix sizes will change only
the rank multiplicities.

We use the following finite approximation form of cone
quasidiagonality: for a separable \(C^*\)-algebra \(R_0\), a finite
subset of \(C_0((0,1])\otimes R_0\), and a positive tolerance,
there is a c.p.c. map to a matrix algebra which is
multiplicative to that tolerance and preserves the norms of
the specified elements to that tolerance.  This is
Voiculescu's cone theorem \cite{Voiculescu}; a direct proof in precisely this
form is \cite[Proposition~33]{Shulman}.

The ideal-covering strategy in the finite cone stage is inspired by
Bosa, Gabe, Sims, and White; compare \cite[Lemma~3.5]{BGSW}.
Their target is \(\mathcal O_\infty\)-stable.  The fixed-cut estimates
and UHF trace-small controls for the present target are proved here,
not imported from that lemma.

\begin{lemma}[A finite cone stage]\label{sm:cone-stage}
Let \(S=C_0((0,1])\otimes P\).  Fix finitely many upper
tests \eqref{sm:upper-tests}, finitely many lower requests
\eqref{sm:error-controllers}, and a finite multiplicativity
test set.  For each lower request one can fix a number
\(h_q>0\), independently of all later stages, such that the
following holds.
For arbitrarily small multiplicativity error there is a
c.p.c. map \(\Phi:S\to M_k(Q)\), with finite constants
\(L_i\) and integers \(M_q\geq1\), satisfying
\begin{equation}\label{sm:stage-bounds}
 \begin{split}
 d_\tau(\Phi(a_i))&\leq L_i d_\tau(b_i),\\
 d_\tau(c_q)&\leq
 M_q d_\tau((\Phi(x_l)-4h_q)_+)
 \qquad(q=(l,\eta)),
 \end{split}
\end{equation}
for every \(\tau\in T(Q)\), with the unnormalized matrix
extension of \(\tau\).
\end{lemma}

\begin{proof}
We first choose the lower threshold, once for each request.
Write
\[
 [c_q]\leq[(z_q-\gamma_q)_+],
 \qquad z_q\in Q(U_{x_l})_+,\quad\gamma_q>0.
\]
For \(z\in X\), denote by \((x_l)_z\) the image in
\(C_0((0,1])\otimes(P/P_z)\), where \(P_z\) is the
primitive ideal labeled by \(z\).  The sets
\[
 O_q(h)=\{z:\|(x_l)_z\|>8h\}
\]
are open and increase to \(U_{x_l}\) as \(h\downarrow0\).
Approximation in this increasing union of ideals allows us
to choose \(h_q>0\) such that
\begin{equation}\label{sm:envelope}
 z'_q=(z_q-\gamma_q/2)_+\in Q(O_q(h_q)),
 \qquad [c_q]\ll[z'_q].
\end{equation}
Keep \(h_q\) fixed from now on.

For a point \(z\in O_q(h_q)\), choose an open neighborhood
\[
 N_z\subseteq O_q(h_q)\cap
 \bigcap_{\{i:(a_i)_z\ne0\}}\supp_X(b_i).
\]
It exists because \((a_i)_z\ne0\) implies \(z\in V_i\),
and \(\supp_X(b_i)\) contains \(V_i\).
Functional calculus on a positive element of \(Q(N_z)\)
nonzero at \(z\) gives positive contractions \(h_z,k_z\)
with
\begin{equation}\label{sm:plateau}
 h_zk_z=k_z,
 \qquad (k_z)_z\ne0,
 \qquad h_z\text{ compactly Cuntz-supported in }Q(N_z).
\end{equation}
For example, first take a nonzero positive cut, make \(h_z\)
equal to one on a smaller spectral region, and support
\(k_z\) inside that region.
The ideal supports of all these \(k_z\)'s cover
\(O_q(h_q)\).  By \eqref{sm:envelope}, finitely many of
them suffice for
\begin{equation}\label{sm:finite-cover}
 [c_q]\leq M_q\sum_{t=1}^{\ell_q}[k_{qt}]
\end{equation}
for some integer \(M_q\).
This uses compact containment in the Cuntz semigroup: the
ideal generated by all the \(k_z\)'s contains \(z'_q\),
and \([c_q]\ll[z'_q]\) is captured by a finite sum with
finite multiplicity.  No compactness of the whole open set
\(O_q(h_q)\) is assumed.

For each selected center \(z\), put
\(S_z=C_0((0,1])\otimes(P/P_z)\), and let
\(\alpha_t\), \(0\leq t\leq1\), be cone scaling:
\[
 (\alpha_t f)(r)=f(tr),
 \qquad \alpha_0=0.
\]
The scaled images of a finite test set form norm-compact
families.  Choose finite nets in those families and their
products.  Cone quasidiagonality, followed by contractivity
to pass from the nets to the families, gives a c.p.c. map
\(\phi_z:S_z\to M_{d_z}\) with the required uniform
multiplicativity error and
\begin{equation}\label{sm:matrix-signal}
 \|\phi_z((x_l)_z)\|>7h_q.
\end{equation}
Define a matrix-valued cone function by
\[
 t\longmapsto\phi_z(\alpha_t(s_z)).
\]
It is norm-continuous and zero at \(t=0\).  Evaluate its
entries at \(h_z\).  The resulting map
\(\psi_z:S\to M_{d_z}(Q)\) is c.p.c.; it has the same
multiplicativity bound because the evaluation map is a
homomorphism.  Its range lies in \(M_{d_z}(\Her(h_z))\).

For the lower estimate, choose a unit vector \(v\in\C^{d_z}\)
which is an eigenvector of the positive matrix
\(\phi_z((x_l)_z)\) for an eigenvalue \(\lambda>7h_q\).
Equation~\eqref{sm:plateau} implies that evaluation on
\(k_z^{1/2}\) is evaluation at parameter one.  Consequently
the rectangular column
\[
 w=(\lambda-4h_q)^{-1/2}\,v\otimes k_z^{1/2}
\]
satisfies
\begin{equation}\label{sm:plateau-witness}
 w^*(\psi_z(x_l)-4h_q)_+w=k_z,
 \qquad \|w\|\leq(3h_q)^{-1/2}.
\end{equation}
Combining these comparisons with
\eqref{sm:finite-cover} gives the second inequality of
\eqref{sm:stage-bounds} for the direct sum of all the
selected blocks.

For an upper test, a block is zero if \((a_i)_z=0\).
Otherwise \(h_z\) is compactly Cuntz-supported in an ideal
contained in \(\Ideal_Q(b_i)\).  It is therefore below
finitely many copies of \(b_i\).  Since the block range
is supported by \(1_{d_z}\otimes h_z\), its rank is at
most \(d_zd_\tau(h_z)\).  Summing these finite bounds gives
the first inequality of \eqref{sm:stage-bounds}.
The direct sum is contractive, since its norm is the maximum
of the block norms.  Its internal matrix trace is the sum
of the block traces.  This proves the lemma.
\end{proof}

The finite stage now has a lower spectral cut independent of the
accuracy.  To turn these maps into a locally small homomorphism, we
choose their UHF fractions next and delay the coordinates until the
prescribed errors are smaller than those fractions.

\begin{proof}[Proof of \Cref{sm:stabilizers}\textup{(ii)}]
At stage \(j\), apply \Cref{sm:cone-stage} to the first
\(j\) upper tests, lower requests, and an increasing dense
list of algebra tests, with multiplicativity error at most
\(j^{-1}\).  Write the resulting map as \(\Phi_j\),
and its constants as \(L_{ij}\) and \(M_{qj}\).
Choose a positive rational \(t_j\) satisfying
\begin{equation}\label{sm:cone-fraction}
 t_j\leq j^{-1},
 \qquad t_jL_{ij}\leq j^{-1}\quad(i\leq j),
 \qquad \beta_{qj}:=t_j/M_{qj}>0.
\end{equation}
After all the stage maps and fractions have been chosen,
choose increasing coordinate thresholds \(N_j\) with
\begin{equation}\label{sm:cone-delay}
 r_q(n)\leq\beta_{qj}
 \qquad(q\leq j, n\geq N_j).
\end{equation}
On \(N_j\leq n<N_{j+1}\), tensor \(\Phi_j\) with a UHF
projection of trace \(t_j\), and fold its finite matrices
to obtain \(\theta_n:S\to Q\).
These maps are c.p.c. and asymptotically multiplicative.
Their quotient is therefore a homomorphism \(\theta\).

For a fixed upper test, \eqref{sm:stage-bounds} and
\eqref{sm:cone-fraction} give
\eqref{sm:upper-goal} eventually with bound \(j^{-1}\).
Thus \(\theta(S(U))\subseteq J(U)\).
For a lower request \(q=(l,\eta)\),
\begin{equation}\label{sm:cone-coordinate-comparison}
 d_\tau(A_{q,n})
 \leq r_q(n)d_\tau(c_q)
 \leq d_\tau((\theta_n(x_l)-4h_q)_+)
\end{equation}
eventually at every coordinate.  Pointwise trace order in
\(\Cu(Q)\) converts this into Cuntz comparison.
Put \(B_{q,n}=(\theta_n(x_l)-2h_q)_+\), choose comparison
witnesses with error at most \(n^{-1}\), and apply
\eqref{sm:witness-bound} with tolerance \(2h_q\).
The resulting witnesses have squared norms at most
\[
 \frac{\sup_n\|A_{q,n}\|+1}{2h_q},
\]
independently of the stage, its centers, and its matrix size.
They prove
\[
 [(|e_l|-\eta)_+]
 \leq[(\theta(x_l)-2h_q)_+]
 \quad\text{in }\Cu(D).
\]
This establishes \eqref{sm:fixed-lower-cut} and then
\eqref{sm:local-and-lower}.  A very small \(\beta_{qj}\)
changes only the delay \(N_j\); it never changes the fixed
cut \(h_q\).
All comparisons may be made in the common block sequence
algebra before folding.  Include the finitely many prescribed
error representatives in each stage's folding tests in
\Cref{st:folding}.  The coordinate conjugacies then preserve
their original quotient classes, and the norm bounds above
are unchanged.  This proves the separate-corner assertion.
\end{proof}

\subsection{Coefficient families}

\begin{corollary}[Simultaneous coefficient domination]
\label{sm:coefficient-domination}
Suppose that one of the hypotheses of
\Cref{sm:stabilizers} holds.  Let \(E_j\) be countably
many bounded multilinear maps
\[
 E_j:S\times Z_{j,1}\times\cdots\times Z_{j,r_j}\longrightarrow D,
\]
where the parameter spaces are separable normed spaces.
Assume that
\[
 E_j(S(U),Z_{j,1},\ldots,Z_{j,r_j})\subseteq J(U)
 \quad\text{for every open }U.
\]
Then the stabilizer may also be chosen so that, for every
\(x\in S_+\),
\begin{equation}\label{sm:all-coefficients}
 E_j(S(U_x),Z_{j,1},\ldots,Z_{j,r_j})
 \subseteq\Ideal_D(\theta(x)).
\end{equation}
The same conclusion holds for norm-continuous separably
parametrized families of such maps, with local norm bounds,
and for any choice of the localized multilinear input.
\end{corollary}

\begin{proof}
Choose a countable family \(\mathcal B\subseteq S_+\) of
positive cuts which is an ideal basis: for every \(x\in S_+\),
the elements of \(\mathcal B\cap\Ideal_S(x)\) generate
\(\Ideal_S(x)\).  Such a family is obtained from a countable
norm-dense positive set and rational cuts, using the
positive-element perturbation lemma.  In particular,
\begin{equation}\label{sm:projected-cover}
 U_x=\bigcup_{y\in\mathcal B\cap\Ideal_S(x)}U_y.
\end{equation}
For each \(y\in\mathcal B\), prescribe the errors
\(E_j(a,z_1,\ldots,z_{r_j})\), where \(a\) runs through
a countable dense set in \(S(U_y)\), and the other variables
run through countable dense sets in their parameter spaces.
These are countably many requests with controlling element
\(y\), to which \Cref{sm:stabilizers} applies.
Continuity gives the desired containment for every input in
\(S(U_y)\) and every parameter value.

Now fix \(x\).  If \(y\in\mathcal B\cap\Ideal_S(x)\),
the homomorphism property gives
\(\theta(y)\in\Ideal_D(\theta(x))\).
By \eqref{sm:projected-cover}, \(S(U_x)\) is the closed sum
of the ideals \(S(U_y)\) in that equation.  Linearity in
the localized input and closedness of the target ideal give
\eqref{sm:all-coefficients}.  Dense parameter sets and local
norm continuity handle the asserted continuous families.
Notice that the argument uses
\(y\in\Ideal_S(x)\), not merely \(U_y\subseteq U_x\).
Thus it does not presume coarse fullness of the cone
stabilizer.
\end{proof}

\begin{remark}\label{sm:multiplier-coefficients}
Two forms of \Cref{sm:coefficient-domination} will be used.
If \(I\) is a separable error algebra and
\(d:S\to D\) is bounded linear with
\(d(S(U))I\subseteq J(U)\), apply the corollary to
\(E(a,b)=d(a)b\), \(b\in I\), and, when needed, to the
adjoint and opposite-product maps.  This gives the full
projected-support bounds on the hereditary support
generators.

Alternatively, let \(H\subseteq D\) be a given
\(\sigma\)-unital hereditary algebra, possibly nonseparable,
and let \(d:S\to M(H)\) be a bounded linear map whose
localized coefficients lie in \(J(U)\).
Choose a countable approximate unit \((u_m)\) for \(H\)
and prescribe \(u_md(a)u_m\) by the corollary.  For
\(h,k\in H\),
\[
 h d(a)k=\lim_m h u_m d(a)u_m k.
\]
Thus all multiplier coefficients are controlled.  The
countability needed here comes from \(\sigma\)-unitality,
not from separability of \(H\).
\end{remark}
